\section{Introduction}
\label{sec:intro}

% Motivation: why calibrated difficulty (not accuracy %) matters in GxP.
Large language models are increasingly proposed as assistants in
GxP-regulated environments such as drug manufacturing, clinical trials,
and regulatory submissions. In these settings an incorrect compliance
determination carries legal and patient-safety consequences. However, the
dominant way we measure LLM competence is a single accuracy percentage on
a fixed test set, and that measure is poorly suited to this setting.
Accuracy conflates easy and hard items, it gives no per-item difficulty,
and it cannot place a new system on a stable scale for comparison across
time or across interventions.

% Gap.
Educational measurement solved an analogous problem decades ago with
\emph{item response theory} (IRT), which separates item difficulty ($b$)
from examinee ability ($\theta$) on a common latent scale
\citep{lord1980,rasch1960}. IRT is standard for human high-stakes
testing, but it is rarely applied to LLM evaluation. Where it is applied,
the usual approach is to fit IRT to an item set that already exists,
either a repurposed human exam or an established NLP benchmark. That
approach is not available here. In a specialized regulatory domain there
is no large pool of human test-takers to calibrate against, and
crowd-sourced items lack the domain subtlety that makes compliance
questions hard.

% Contributions.
This paper builds a difficulty-calibrated compliance item bank end to
end, uses it as a measuring instrument, and then tests the instrument
itself. The contributions are the following.

\begin{enumerate}[leftmargin=1.4em]
  \item \textbf{Automated item generation with a difficulty mandate.}
  I generate compliance scenarios under an explicit ``Red Herring plus
  subtle violation'' construction (Section~\ref{sec:itemgen}). Every hard
  item must contain a compliant-but-suspicious element designed to
  mislead, and the true violation must be a procedural or data-integrity
  subtlety rather than an obvious lapse. A parallel track converts real
  FDA warning letters into items, which grounds the synthetic bank in
  actual enforcement.

  \item \textbf{IRT calibration with a synthetic examinee population.}
  Lacking human test-takers, I administer each item to a population of
  synthetic AI examinees and estimate difficulty from their response
  pattern (Sections~\ref{sec:calibration}--\ref{sec:pipeline}). I fit a
  Rasch (1PL) model, anchor it to a fixed baseline examinee, and apply
  classical item-fit quality control to freeze a healthy bank of
  \textbf{1{,}284} items.

  \item \textbf{$\theta$ as an optimization metric, with a
  format-versus-reasoning decomposition.} Treating the calibrated bank as
  a fixed instrument, I measure how far two kinds of intervention move a
  model's ability: system-prompt strictness
  (Section~\ref{sec:deltatheta}) and agent architecture, including
  retrieval, self-critique, and step decomposition
  (Section~\ref{sec:phase3b}). Grading the same responses with a strict
  judge and a lenient judge separates the gain that comes from
  \emph{citation format} from the gain that reflects genuine
  \emph{compliance reasoning}.

  \item \textbf{Validation of the instrument.} Because every result rests
  on an LLM judge, I audit that judge against two independent evaluation
  frameworks (Section~\ref{sec:judgeaudit}), and I document the
  human expert-review protocol now collecting labels
  (Section~\ref{sec:expertreview}).
\end{enumerate}

The paper has three parts. Part~1
(Sections~\ref{sec:itemgen}--\ref{sec:results}) builds the
\emph{instrument} and describes what the calibrated bank looks like.
Part~2 (Sections~\ref{sec:deltatheta}--\ref{sec:phase3b}) runs two
\emph{experiments} on that instrument. Part~3
(Sections~\ref{sec:judgeaudit}--\ref{sec:expertreview}) turns the
question around and asks how far the instrument itself can be trusted.
The loop is closed, in that the same model family generates, answers, and
judges the items, and I return to what that costs in
Section~\ref{sec:discussion}.
